\begin{table}[H]
    \centering
    \caption{Rerata Korelasi \textit{Off-Diagonal} Antarsubtes secara Kumulatif Seiring Bertambahnya Tahap Pengumpulan Data}
    \label{tab:saturasi_tahap}
    \begin{spacing}{1.0}
    \resizebox{\textwidth}{!}{%
    \begin{tabular}{cllcc}
        \toprule
        Tahap & Nama Tahap & Kategori & $N$ Kumulatif & $\bar{r}$ \textit{off-diagonal} \\
        \midrule
        1 & Utas Tunggal \& DM Pribadi     & Campuran        & 216 & 0,297 \\
        2 & Utas Khusus Unpad              & PTN Khusus      & 249 & 0,300 \\
        3 & Utas Khusus PENS               & PTN Khusus      & 255 & 0,298 \\
        4 & Utas Khusus Unsoed             & PTN Khusus      & 289 & 0,298 \\
        5 & Utas Khusus Unesa              & PTN Khusus      & 300 & 0,293 \\
        6 & Utas Khusus Unnes              & PTN Khusus      & 320 & 0,286 \\
        7 & Utas Khusus ITS                & PTN Khusus      & 416 & 0,345 \\
        8 & Utas Khusus Undip              & PTN Khusus      & 567 & 0,316 \\
        9 & Utas Khusus IPB                & PTN Khusus      & 613 & 0,304 \\
        10 & Utas Khusus UNS                & PTN Khusus      & 675 & 0,294 \\
        11 & Utas Khusus UB                 & PTN Khusus      & 816 & 0,269 \\
        12 & Utas Khusus UPI                & PTN Khusus      & 847 & 0,267 \\
        13 & Utas Khusus UGM                & PTN Khusus      & 992 & 0,280 \\
        14 & Utas Khusus Unair              & PTN Khusus      & 1.095 & 0,270 \\
        15 & Utas Khusus PNJ                & PTN Khusus      & 1.107 & 0,269 \\
        16 & Utas Khusus UI                 & PTN Khusus      & 1.319 & 0,286 \\
        17 & Utas Khusus ITB                & PTN Khusus      & 1.459 & 0,305 \\
        18 & Utas Terbuka Inklusif (@sbmptnfess) & Inklusif        & 1.709 & 0,309 \\
        19 & Akun Publik @utbkcak           & Inklusif        & 1.886 & 0,300 \\
        \bottomrule
    \end{tabular}%
    }
    \end{spacing}
    \vspace{2mm}
    \begin{flushleft}
    \footnotesize\textit{Catatan.} $N$ kumulatif adalah jumlah entri dari Tahap 1 sampai tahap terkait (tahap disusun berurutan sesuai kronologi pengumpulan). Nilai $\bar{r}$ \textit{off-diagonal}, yaitu rerata korelasi di luar diagonal matriks korelasi \textit{raw sample}, dihitung ulang dari entri kumulatif itu.
    \end{flushleft}
\end{table}
